\chapter{Proposed Temporal-State-Aware Compact Routing Method}

\section{Introduction}

This chapter presents the proposed temporal-state-aware compact routing method for quantum-native entanglement overlay networks. The method extends compact quantum routing by incorporating the time-dependent state of stored artificial e-links into route selection and overlay-maintenance decisions.

The proposed framework is motivated by the observation that an artificial e-link may remain structurally present in the routing table while its effective routing utility deteriorates because of aging, decoherence, fidelity degradation, memory waiting time, and delayed consumption. Therefore, a routing decision based only on topological reachability or static entangling cost may select a route that is unlikely to satisfy the required end-to-end fidelity at completion time.

The proposed method addresses this limitation through five coordinated operations:

\begin{enumerate}
    \item monitoring the temporal state of each stored artificial e-link;
    \item constructing compact-routing candidate paths;
    \item filtering candidates using temporal feasibility conditions;
    \item selecting the minimum-cost temporally feasible route; and
    \item regenerating or replacing critical artificial e-links when necessary.
\end{enumerate}

The method operates on the artificial entanglement overlay while using physical-network information to estimate generation delay, regeneration cost, fidelity evolution, and memory consumption.

At each time slot, the controller updates the state of every artificial e-link according to

\begin{equation}
    \mathcal{L}_{ij}(t)
    =
    \left[
    x_{ij}(t),
    A_{ij}(t),
    F_{ij}^{\mathrm{A},0},
    F_{ij}^{\mathrm{eff}}(t),
    T_{2,ij},
    B_{ij}(t),
    D_{ij}^{\mathrm{regen}},
    C_{ij}^{\mathrm{regen}},
    I_{ij}^{\mathrm{struct}}
    \right].
\end{equation}

Based on this temporal state, the controller classifies each artificial e-link as

\begin{equation}
    \mathrm{State}_{ij}(t)
    \in
    \{
    \mathrm{Safe},
    \mathrm{Warning},
    \mathrm{Critical},
    \mathrm{Expired}
    \}.
\end{equation}

Only safe or admissible warning-state links are directly considered during route construction. Critical links may trigger route repair or proactive regeneration, whereas expired links are removed from the temporally usable overlay.

The proposed routing process is expressed as

\begin{equation}
    \mathrm{Monitor}
    \rightarrow
    \mathrm{Construct}
    \rightarrow
    \mathrm{Filter}
    \rightarrow
    \mathrm{Select}
    \rightarrow
    \mathrm{Maintain}.
\end{equation}

The method preserves the compact-routing principle by limiting route discovery to local e-neighbors, anchor nodes, and tracked long-range artificial links. It does not require every ESP to maintain complete global routing information.

The main design challenge is therefore to improve temporal reliability without converting the compact-routing architecture into a fully centralized global-state routing system.

This chapter first defines the design objectives and overall framework. It then describes temporal-state monitoring, candidate-route construction, temporal-feasibility filtering, time-dependent route selection, proactive regeneration, overlay maintenance, algorithmic steps, and computational complexity.
\section{Design Objectives}

The proposed method is designed to improve the reliability of compact quantum routing by incorporating temporal information of artificial e-links while preserving the scalability advantages of compact routing. The method does not require every Entanglement Service Provider to maintain complete global knowledge of the network state.

The main design objectives are described as follows.

\subsection{Temporal Reliability}

The first objective is to prevent the routing protocol from selecting artificial e-links that are structurally available but expected to become unusable before route completion.

For every selected route $R_{s,d}(t)$, the proposed method aims to satisfy the coherence-budget constraint

\begin{equation}
    B_{ij}(t)
    \geq
    \delta_{\mathrm{safe}},
    \qquad
    \forall (i,j)
    \in
    R_{s,d}(t),
\end{equation}

where $B_{ij}(t)$ denotes the remaining coherence budget of artificial e-link $(i,j)$ and $\delta_{\mathrm{safe}}$ represents the required safety margin.

In addition, the selected route must satisfy the predicted end-to-end fidelity requirement at the expected completion time:

\begin{equation}
    \widehat{F}_{s,d}^{\mathrm{end}}
    \left(
    t+
    \widehat{D}_{s,d}^{\mathrm{route}}(t)
    \right)
    \geq
    F_{\min}^{\mathrm{end}}.
\end{equation}

These constraints ensure that routing decisions consider future link usability rather than only the current availability state.

\subsection{Fidelity Preservation}

The second objective is to preserve the quality of distributed entanglement by considering both the current fidelity and the predicted fidelity at the expected consumption time.

The routing method therefore avoids selecting an artificial e-link solely because its instantaneous fidelity satisfies the threshold. Instead, it evaluates whether the link quality remains sufficient after additional waiting, coordination, purification, and swapping operations.

The fidelity-preservation objective can be expressed as

\begin{equation}
    \max
    \widehat{F}_{s,d}^{\mathrm{end}}
    \left(
    t+
    \widehat{D}_{s,d}^{\mathrm{route}}(t)
    \right).
\end{equation}

This objective is incorporated into the proposed time-dependent entangling-cost metric rather than optimized independently. Therefore, high-fidelity and temporally stable routes receive lower routing costs during route selection.

\subsection{Compact Routing-State Preservation}

The proposed method must preserve the compact-routing property of the underlying quantum-native architecture.

Let $\mathcal{T}_i(t)$ denote the routing table maintained by ESP $v_i$. The routing-state constraint is defined as

\begin{equation}
    |\mathcal{T}_i(t)|
    \leq
    c_{\mathrm{tab}}
    \sqrt{n_e}\log n_e,
\end{equation}

where $n_e=|V_2|$ represents the number of ESPs.

The method therefore relies on local e-neighbors, anchor information, and selected tracked e-links instead of maintaining complete all-pairs routing information.

\subsection{Low Route-Selection Cost}

The method aims to select a temporally feasible route with minimum time-dependent entangling cost.

For each demand $(s,d)$, the route-selection objective is

\begin{equation}
    R_{s,d}^{*}(t)
    =
    \arg\min_{
    R
    \in
    \mathcal{R}_{s,d}^{\mathrm{TF}}(t)
    }
    w_t(R).
\end{equation}

The metric $w_t(R)$ considers fidelity degradation, e-link age, coherence-budget risk, regeneration cost, and structural routing cost.

\subsection{Controlled Regeneration Overhead}

Frequent regeneration may improve temporal freshness but can consume quantum-memory capacity, elementary-link generation resources, classical coordination time, and Bell-state measurement resources.

The proposed method therefore limits the number of regeneration operations in each time slot:

\begin{equation}
    \sum_{
    (i,j)
    \in
    E_{\mathrm{A}}(t)
    }
    r_{ij}(t)
    \leq
    R_{\max}.
\end{equation}

Regeneration is prioritized for artificial e-links with high temporal risk and high structural importance.

\subsection{Reduced Request Rejection}

The routing method aims to reduce entanglement-request rejection caused by stale links, insufficient coherence budget, or fidelity-threshold violations.

Let $N_{\mathrm{reject}}(t)$ denote the number of rejected requests at time $t$. The long-term request-rejection ratio is

\begin{equation}
    P_{\mathrm{reject}}
    =
    \frac{
    \sum_{t=1}^{T}
    N_{\mathrm{reject}}(t)
    }{
    \sum_{t=1}^{T}
    |\mathcal{D}(t)|
    }.
\end{equation}

The proposed method seeks to minimize this ratio while respecting memory and regeneration constraints.

\subsection{Robustness to Dynamic E-Link Conditions}

The artificial overlay may change because of successful generation, consumption, expiration, memory failure, or regeneration.

The method adapts route selection whenever the state of a relevant e-link changes. However, unnecessary route recomputation should be avoided.

A route update is triggered only when at least one of the following conditions occurs:

\begin{itemize}
    \item a selected e-link becomes critical or expired;
    \item the predicted end-to-end fidelity falls below the required threshold;
    \item a lower-cost temporally feasible route becomes available;
    \item an anchor or tracked e-link is regenerated; or
    \item the current route cannot satisfy the completion-delay constraint.
\end{itemize}

\subsection{Scalable Decision Making}

The proposed method should remain computationally practical as the number of ESPs increases.

Instead of solving the complete joint optimization problem at every time slot, the framework separates the decision process into

\begin{equation}
    \mathrm{State\ Update}
    \rightarrow
    \mathrm{Candidate\ Construction}
    \rightarrow
    \mathrm{Feasibility\ Filtering}
    \rightarrow
    \mathrm{Cost\ Based\ Selection}.
\end{equation}

This decomposition reduces the number of candidate routes requiring detailed temporal evaluation.

The overall design objective can be summarized as

\begin{equation}
    \min
    \left[
    \lambda_1
    \overline{w}
    +
    \lambda_2
    P_{\mathrm{reject}}
    +
    \lambda_3
    C_{\mathrm{regen}}
    +
    \lambda_4
    \overline{ES}
    \right],
\end{equation}

subject to

\begin{align}
    & \widehat{F}_{s,d}^{\mathrm{end}}
    \geq
    F_{\min}^{\mathrm{end}},\\
    & B_{ij}(t)
    \geq
    \delta_{\mathrm{safe}},\\
    & m_i(t)
    \leq
    M_i,\\
    & |\mathcal{T}_i(t)|
    \leq
    c_{\mathrm{tab}}
    \sqrt{n_e}\log n_e,\\
    & \sum_{(i,j)}r_{ij}(t)
    \leq
    R_{\max}.
\end{align}

Here, $\overline{w}$ is the average selected-route cost, $C_{\mathrm{regen}}$ is the total regeneration overhead, and $\overline{ES}$ is the average temporal entangling stretch.
\section{Overview of the Proposed Framework}

The proposed framework integrates compact routing, temporal e-link monitoring, fidelity prediction, coherence-budget evaluation, and selective overlay maintenance.

The framework operates in discrete time slots. At the beginning of each time slot, the controller updates the state of every stored artificial e-link and constructs the temporally usable overlay.

The complete processing sequence is represented as

\begin{equation}
    G_{\mathrm{A}}(t)
    \rightarrow
    \mathcal{L}(t)
    \rightarrow
    G_{\mathrm{U}}(t)
    \rightarrow
    \mathcal{R}_{s,d}^{\mathrm{cand}}(t)
    \rightarrow
    \mathcal{R}_{s,d}^{\mathrm{TF}}(t)
    \rightarrow
    R_{s,d}^{*}(t).
\end{equation}

Here,

\begin{itemize}
    \item $G_{\mathrm{A}}(t)$ is the currently stored artificial overlay;
    \item $\mathcal{L}(t)$ is the collection of temporal states of artificial e-links;
    \item $G_{\mathrm{U}}(t)$ is the temporally usable overlay;
    \item $\mathcal{R}_{s,d}^{\mathrm{cand}}(t)$ is the candidate-route set;
    \item $\mathcal{R}_{s,d}^{\mathrm{TF}}(t)$ is the temporally feasible route set; and
    \item $R_{s,d}^{*}(t)$ is the selected minimum-cost route.
\end{itemize}

The framework contains two interacting control processes:

\begin{itemize}
    \item Reactive routing;
    \item Proactive overlay maintenance.
\end{itemize}

Reactive routing is executed when an entanglement request arrives. It constructs candidate compact routes, evaluates their temporal conditions, removes infeasible routes, and selects the lowest-cost temporally feasible route.

Proactive overlay maintenance is executed periodically or when a critical state change occurs. It identifies artificial e-links that should be regenerated, replaced, or discarded before they cause routing failure.

The framework receives the following inputs:

\begin{equation}
    \mathcal{I}(t)
    =
    \left[
    G_{\mathrm{P}},
    G_{\mathrm{A}}(t),
    \mathcal{D}(t),
    \mathcal{T}(t),
    \mathcal{M}(t),
    \mathcal{Q}(t)
    \right],
\end{equation}

where

\begin{itemize}
    \item $G_{\mathrm{P}}$ is the physical quantum network;
    \item $G_{\mathrm{A}}(t)$ is the current artificial overlay;
    \item $\mathcal{D}(t)$ is the set of active entanglement requests;
    \item $\mathcal{T}(t)$ is the compact-routing-table state;
    \item $\mathcal{M}(t)$ is the quantum-memory state; and
    \item $\mathcal{Q}(t)$ is the observed or predicted network-quality state.
\end{itemize}

The output of the framework is

\begin{equation}
    \mathcal{O}(t)
    =
    \left[
    \mathcal{Y}(t),
    \mathcal{R}^{\mathrm{regen}}(t),
    \mathcal{E}^{\mathrm{drop}}(t),
    \mathcal{T}(t+1)
    \right],
\end{equation}

where

\begin{itemize}
    \item $\mathcal{Y}(t)$ is the set of selected routes;
    \item $\mathcal{R}^{\mathrm{regen}}(t)$ is the regeneration-action set;
    \item $\mathcal{E}^{\mathrm{drop}}(t)$ is the set of discarded e-links; and
    \item $\mathcal{T}(t+1)$ is the updated compact-routing-table state.
\end{itemize}

The proposed framework is divided into six functional modules.

\subsection{Temporal-State Monitoring Module}

This module updates the age, effective fidelity, predicted consumption fidelity, remaining coherence budget, and stale-link state of every stored artificial e-link.

Its output is the temporally usable artificial overlay

\begin{equation}
    G_{\mathrm{U}}(t)
    =
    \left(
    V_2,
    E_{\mathrm{U}}(t)
    \right).
\end{equation}

\subsection{Compact Candidate-Construction Module}

This module constructs a limited number of candidate routes using only the routing information permitted by the Partial-Anchor or Full-Anchor compact-routing scheme.
\subsection{Compact Candidate-Construction Module}

The candidate-route set generated by this module is represented as

\begin{equation}
    \mathcal{R}_{s,d}^{\mathrm{cand}}(t).
\end{equation}

This restriction preserves compact-routing state requirements and avoids unrestricted global path enumeration.

\subsection{Temporal-Feasibility Module}

This module evaluates each candidate route according to the temporal route-feasibility indicator

\begin{equation}
    \chi_{s,d}
    \left(
    R_{s,d}(t)
    \right).
\end{equation}

Routes that violate fidelity, coherence-budget, memory, delay, or e-link availability constraints are removed.

The remaining route set is

\begin{equation}
    \mathcal{R}_{s,d}^{\mathrm{TF}}(t)
    \subseteq
    \mathcal{R}_{s,d}^{\mathrm{cand}}(t).
\end{equation}

\subsection{Time-Dependent Route-Selection Module}

For each temporally feasible route, this module calculates the accumulated time-dependent entangling cost:

\begin{equation}
    w_t(R)
    =
    \sum_{(i,j)\in R}
    w_t(i,j).
\end{equation}

The selected route is

\begin{equation}
    R_{s,d}^{*}(t)
    =
    \arg\min_{
    R
    \in
    \mathcal{R}_{s,d}^{\mathrm{TF}}(t)
    }
    w_t(R).
\end{equation}

\subsection{Route-Repair Module}

When no directly feasible route exists, this module identifies routes that can become feasible through limited regeneration or purification operations.

The route-repair process is performed only when

\begin{equation}
    \widehat{D}_{s,d}^{\mathrm{repair}}(t)
    \leq
    D_{\max}^{\mathrm{repair}}
\end{equation}

and

\begin{equation}
    C_{s,d}^{\mathrm{repair}}(t)
    \leq
    C_{\max}^{\mathrm{repair}}.
\end{equation}

Otherwise, the request is delayed or rejected.

\subsection{Overlay-Maintenance Module}

This module assigns regeneration priority to critical artificial e-links according to temporal risk and structural importance.

The maintenance priority is defined as

\begin{equation}
    \Pi_{ij}(t)
    =
    \beta_1
    C_B(i,j,t)
    +
    \beta_2
    C_F(i,j,t)
    +
    \beta_3
    I_{ij}^{\mathrm{struct}}
    +
    \beta_4
    H_{ij}^{\mathrm{use}}(t),
\end{equation}

where $H_{ij}^{\mathrm{use}}(t)$ represents the recent or predicted utilization level of the artificial e-link.

The weights satisfy

\begin{equation}
    \sum_{k=1}^{4}
    \beta_k
    =
    1.
\end{equation}

The highest-priority links are regenerated subject to quantum-memory availability, generation-attempt budget, and control-overhead limitations.

The framework is intentionally modular. Different fidelity models, delay predictors, compact-routing structures, or regeneration policies can be integrated without changing the overall processing sequence.

The proposed method is a model-based routing and control heuristic. It does not require a trained machine-learning model. Any prediction component used in the implementation is derived from measured averages, analytical estimates, or simulation parameters.
\section{Temporal E-Link State Monitoring}

The temporal e-link state-monitoring module continuously evaluates the condition of every artificial e-link stored in the quantum memories of the Entanglement Service Providers.

Its purpose is to distinguish between an artificial e-link that is merely present in the overlay and an e-link that remains sufficiently reliable for future routing and entanglement-swapping operations.

Let the set of currently stored artificial e-links be

\begin{equation}
    E_{\mathrm{A}}(t).
\end{equation}

For each artificial e-link $(i,j)\in E_{\mathrm{A}}(t)$, the controller maintains the temporal state record

\begin{equation}
    \mathcal{L}_{ij}(t)
    =
    \left[
    x_{ij}(t),
    A_{ij}(t),
    F_{ij}^{\mathrm{A},0},
    F_{ij}^{\mathrm{eff}}(t),
    T_{2,ij},
    \widehat{F}_{ij}^{\mathrm{cons}}(t),
    B_{ij}(t),
    H_{ij}^{\mathrm{use}}(t),
    S_{ij}^{\mathrm{state}}(t)
    \right],
\end{equation}

where

\begin{itemize}
    \item $x_{ij}(t)$ is the current storage and availability indicator;
    \item $A_{ij}(t)$ is the current e-link age;
    \item $F_{ij}^{\mathrm{A},0}$ is the initial fidelity after e-link generation;
    \item $F_{ij}^{\mathrm{eff}}(t)$ is the current effective fidelity;
    \item $T_{2,ij}$ is the coherence time of the stored entanglement;
    \item $\widehat{F}_{ij}^{\mathrm{cons}}(t)$ is the predicted fidelity at consumption time;
    \item $B_{ij}(t)$ is the remaining coherence budget;
    \item $H_{ij}^{\mathrm{use}}(t)$ is the recent or predicted utilization score; and
    \item $S_{ij}^{\mathrm{state}}(t)$ is the temporal-state classification.
\end{itemize}

At the end of each time slot, the age of an active e-link is updated as

\begin{equation}
    A_{ij}(t+1)
    =
    A_{ij}(t)
    +
    1,
\end{equation}

provided that the link remains stored and is neither consumed nor regenerated.

When regeneration succeeds, the age is reset:

\begin{equation}
    A_{ij}(t+1)
    =
    0.
\end{equation}

The current effective fidelity is estimated using the selected decoherence model:

\begin{equation}
    F_{ij}^{\mathrm{eff}}(t)
    =
    f_{\mathrm{dec}}
    \left(
    F_{ij}^{\mathrm{A},0},
    A_{ij}(t),
    T_{2,ij}
    \right).
\end{equation}

The controller also estimates the additional delay before the link is expected to be consumed:

\begin{equation}
    \widehat{D}_{ij}^{\mathrm{use}}(t).
\end{equation}

This delay may be estimated from recent queueing observations, route depth, expected swapping order, and classical coordination delay.

The predicted fidelity at consumption time is calculated as

\begin{equation}
    \widehat{F}_{ij}^{\mathrm{cons}}(t)
    =
    f_{\mathrm{dec}}
    \left(
    F_{ij}^{\mathrm{eff}}(t),
    \widehat{D}_{ij}^{\mathrm{use}}(t),
    T_{2,ij}
    \right).
\end{equation}

The remaining coherence budget is updated according to

\begin{equation}
    B_{ij}(t)
    =
    T_{2,ij}
    -
    A_{ij}(t)\Delta
    -
    \widehat{D}_{ij}^{\mathrm{use}}(t).
\end{equation}

\subsection{Temporal-State Classification and Link Availability}

To reduce sensitivity to short-term fluctuations, the utilization score is updated using an exponential moving average:

\begin{equation}
    H_{ij}^{\mathrm{use}}(t)
    =
    \gamma_H
    H_{ij}^{\mathrm{use}}(t-1)
    +
    \left(
    1-\gamma_H
    \right)
    I_{ij}^{\mathrm{used}}(t),
\end{equation}

where

\begin{equation}
    0
    \leq
    \gamma_H
    <
    1
\end{equation}

and

\begin{equation}
    I_{ij}^{\mathrm{used}}(t)
    =
    \begin{cases}
        1, & \text{if the e-link is selected or consumed at time $t$},\\
        0, & \text{otherwise}.
    \end{cases}
\end{equation}

The temporal state of an artificial e-link is classified according to predicted consumption fidelity and remaining coherence budget.

A link is classified as safe when

\begin{equation}
    B_{ij}(t)
    \geq
    B_{\mathrm{high}}
\end{equation}

and

\begin{equation}
    \widehat{F}_{ij}^{\mathrm{cons}}(t)
    \geq
    F_{\mathrm{high}}.
\end{equation}

A link is classified as warning when it remains usable but approaches temporal degradation:

\begin{equation}
    B_{\mathrm{low}}
    \leq
    B_{ij}(t)
    <
    B_{\mathrm{high}}
\end{equation}

or

\begin{equation}
    F_{\min}^{\mathrm{link}}
    \leq
    \widehat{F}_{ij}^{\mathrm{cons}}(t)
    <
    F_{\mathrm{high}}.
\end{equation}

A link is classified as critical when it is still potentially recoverable but has a high risk of becoming unusable:

\begin{equation}
    0
    <
    B_{ij}(t)
    <
    B_{\mathrm{low}}
\end{equation}

or

\begin{equation}
    \widehat{F}_{ij}^{\mathrm{cons}}(t)
    \approx
    F_{\min}^{\mathrm{link}}.
\end{equation}

Critical links are not directly preferred for route selection and may trigger proactive regeneration or purification.

A link is classified as expired when

\begin{equation}
    B_{ij}(t)
    \leq
    0
\end{equation}

or

\begin{equation}
    x_{ij}(t)
    =
    0.
\end{equation}

The resulting temporal-state classification is

\begin{equation}
    S_{ij}^{\mathrm{state}}(t)
    =
    \begin{cases}
        \mathrm{Safe},
        & \text{if high-quality conditions are satisfied},\\
        \mathrm{Warning},
        & \text{if the link remains usable with degradation risk},\\
        \mathrm{Critical},
        & \text{if recovery actions may restore usability},\\
        \mathrm{Expired},
        & \text{if the link is unusable}.
    \end{cases}
\end{equation}

The direct-use indicator is defined as

\begin{equation}
    u_{ij}^{\mathrm{direct}}(t)
    =
    \begin{cases}
        1,
        &
        S_{ij}^{\mathrm{state}}(t)
        \in
        \{\mathrm{Safe},\mathrm{Warning}\},\\
        0,
        &
        \text{otherwise}.
    \end{cases}
\end{equation}

The regeneration-candidate indicator is defined as

\begin{equation}
    u_{ij}^{\mathrm{regen}}(t)
    =
    \begin{cases}
        1,
        &
        S_{ij}^{\mathrm{state}}(t)
        =
        \mathrm{Critical},\\
        0,
        &
        \text{otherwise}.
    \end{cases}
\end{equation}

Expired links are removed from route construction and may be discarded from memory:

\begin{equation}
    u_{ij}^{\mathrm{drop}}(t)
    =
    \mathbf{1}
    \left[
    S_{ij}^{\mathrm{state}}(t)
    =
    \mathrm{Expired}
    \right].
\end{equation}

The temporally usable edge set is therefore

\begin{equation}
    E_{\mathrm{U}}(t)
    =
    \left\{
    (i,j)
    \in
    E_{\mathrm{A}}(t)
    :
    u_{ij}^{\mathrm{direct}}(t)=1
    \right\}.
\end{equation}

The temporal-state monitor operates periodically. Let

\begin{equation}
    \tau_{\mathrm{mon}}
\end{equation}

denote the monitoring interval.

A shorter monitoring interval improves responsiveness but increases control and state-update overhead. A longer interval reduces overhead but may allow rapidly degrading links to remain incorrectly classified.

Therefore, the monitoring interval is evaluated experimentally under different coherence times, request rates, and overlay sizes.

The output of this module is a synchronized temporal-state table used by candidate-route construction, temporal-feasibility filtering, route selection, and proactive regeneration.
\section{Candidate Route Construction}

The candidate-route construction module generates a limited set of compact-routing paths between a source ESP and a destination ESP. Its purpose is to avoid unrestricted enumeration of all possible paths while preserving the structural properties of the underlying Partial-Anchor or Full-Anchor routing scheme.

For an entanglement request between source ESP $v_s$ and destination ESP $v_d$, the candidate-route set at time $t$ is denoted by

\begin{equation}
    \mathcal{R}_{s,d}^{\mathrm{cand}}(t).
\end{equation}

Candidate routes are constructed over the currently available compact-overlay information. The module uses the temporally usable overlay

\begin{equation}
    G_{\mathrm{U}}(t)
    =
    \left(
    V_2,
    E_{\mathrm{U}}(t)
    \right),
\end{equation}

where $E_{\mathrm{U}}(t)$ contains artificial e-links classified as safe or warning.

The candidate-generation process first checks whether the source and destination are directly connected by a usable artificial e-link. If

\begin{equation}
    (v_s,v_d)
    \in
    E_{\mathrm{U}}(t),
\end{equation}

then the direct route

\begin{equation}
    R_{s,d}^{\mathrm{dir}}(t)
    =
    \left(
    v_s,
    v_d
    \right)
\end{equation}

is inserted into the candidate set.

If no direct usable e-link exists, the controller searches through local e-neighbors, anchor ESPs, and tracked long-range links according to the selected compact-routing architecture.

Let

\begin{equation}
    \mathcal{N}_{\mathrm{U}}(v_i,t)
\end{equation}

denote the usable e-neighborhood of ESP $v_i$:

\begin{equation}
    \mathcal{N}_{\mathrm{U}}(v_i,t)
    =
    \left\{
    v_j
    \in
    V_2
    :
    (v_i,v_j)
    \in
    E_{\mathrm{U}}(t)
    \right\}.
\end{equation}

A one-intermediate-node candidate route has the form

\begin{equation}
    R_{s,d}^{(1)}(t)
    =
    \left(
    v_s,
    v_k,
    v_d
    \right),
\end{equation}

where

\begin{equation}
    v_k
    \in
    \mathcal{N}_{\mathrm{U}}(v_s,t)
    \cap
    \mathcal{N}_{\mathrm{U}}(v_d,t).
\end{equation}

For longer compact routes, the controller follows the routing structure of the selected overlay scheme. In the Partial-Anchor scheme, route construction is restricted by local e-neighborhood information and selected anchor links. In the Full-Anchor scheme, additional tracked long-range links are considered during candidate generation.

\subsection{Candidate Construction in the Partial-Anchor Scheme}

In the Partial-Anchor Scheme, each ESP maintains local e-neighbor information, while a subset of ESPs acts as anchors.

Let

\begin{equation}
    \mathcal{T}
    \subseteq
    V_2
\end{equation}

denote the anchor set.

The usable source-anchor set is

\begin{equation}
    \mathcal{T}_{s}^{\mathrm{U}}(t)
    =
    \left\{
    v_a
    \in
    \mathcal{T}
    :
    v_a
    \in
    \mathcal{N}_{\mathrm{U}}(v_s,t)
    \right\}.
\end{equation}

Similarly, the usable destination-anchor set is

\begin{equation}
    \mathcal{T}_{d}^{\mathrm{U}}(t)
    =
    \left\{
    v_b
    \in
    \mathcal{T}
    :
    v_b
    \in
    \mathcal{N}_{\mathrm{U}}(v_d,t)
    \right\}.
\end{equation}

A Partial-Anchor candidate route can be written as

\begin{equation}
    R_{s,d}^{\mathrm{PA}}(t)
    =
    \left(
    v_s,
    v_a,
    \ldots,
    v_b,
    v_d
    \right),
\end{equation}

where

\begin{equation}
    v_a
    \in
    \mathcal{T}_{s}^{\mathrm{U}}(t)
\end{equation}

and

\begin{equation}
    v_b
    \in
    \mathcal{T}_{d}^{\mathrm{U}}(t).
\end{equation}

The internal anchor segment is obtained from the currently usable anchor overlay.

Let

\begin{equation}
    G_{\mathrm{anchor}}^{\mathrm{U}}(t)
\end{equation}

denote the usable subgraph formed by anchor ESPs and usable anchor e-links.

A bounded number of low-cost anchor paths is generated between $v_a$ and $v_b$.

\subsection{Candidate Construction in the Full-Anchor Scheme}

In the Full-Anchor Scheme, every ESP may maintain local e-links and selected tracked long-range e-links.

Let

\begin{equation}
    E_{\mathrm{tracked}}^{\mathrm{U}}(t)
\end{equation}

denote the set of currently usable tracked long-range artificial e-links.

The Full-Anchor candidate set is constructed using

\begin{equation}
    E_{\mathrm{local}}^{\mathrm{U}}(t)
    \cup
    E_{\mathrm{tracked}}^{\mathrm{U}}(t).
\end{equation}

A candidate route is constructed by combining local source segments, tracked long-range segments, and local destination segments.

A general candidate route is represented by

\begin{equation}
    R_{s,d}^{\mathrm{FA}}(t)
    =
    \left(
    v_s,
    v_{r_1},
    \ldots,
    v_{r_h},
    v_d
    \right),
\end{equation}

where every consecutive node pair corresponds to a usable local or tracked artificial e-link.
\subsection{Bounded Candidate-Set Size}

To control computational overhead, the controller retains at most

\begin{equation}
    K_{\max}
\end{equation}

candidate routes for each source--destination request.

Therefore,

\begin{equation}
    \left|
    \mathcal{R}_{s,d}^{\mathrm{cand}}(t)
    \right|
    \leq
    K_{\max}.
\end{equation}

Candidate routes are initially ranked using a lightweight preliminary score

\begin{equation}
    \widetilde{w}_t(R)
    =
    \sum_{(i,j)\in R}
    \widetilde{w}_t(i,j),
\end{equation}

where $\widetilde{w}_t(i,j)$ includes only readily available local information such as hop count, current fidelity, normalized age, and direct-use state.

A possible preliminary link score is

\begin{equation}
    \widetilde{w}_t(i,j)
    =
    \eta_h
    +
    \eta_F
    \left(
    1-F_{ij}^{\mathrm{eff}}(t)
    \right)
    +
    \eta_A
    \frac{
    A_{ij}(t)
    }{
    A_{ij}^{\max}
    }.
\end{equation}

The coefficients satisfy

\begin{equation}
    \eta_h+\eta_F+\eta_A=1.
\end{equation}

This preliminary score is used only for candidate pruning. Final route selection uses the complete time-dependent entangling-cost metric.

To avoid routing loops, every candidate route must contain unique nodes:

\begin{equation}
    v_m
    \neq
    v_n,
    \qquad
    \forall m\neq n .
\end{equation}

A maximum overlay-hop constraint is also applied:

\begin{equation}
    |R|-1
    \leq
    H_{\max},
\end{equation}

where $H_{\max}$ is the maximum number of artificial e-links allowed in a candidate route.

The final candidate set is

\begin{equation}
    \mathcal{R}_{s,d}^{\mathrm{cand}}(t)
    =
    \operatorname{TopK}_{K_{\max}}
    \left[
    \widetilde{w}_t(R)
    \right],
\end{equation}

subject to compact-routing, loop-free, direct-use, and maximum-hop constraints.

The output of this module is a small set of structurally valid candidate routes. These routes are passed to the temporal-feasibility module, where predicted consumption fidelity, remaining coherence budget, delay, memory availability, and end-to-end fidelity are evaluated in detail.

\subsection{Availability Test}

The first test verifies that every required artificial e-link is currently available for direct use:

\begin{equation}
    u_{ij}^{\mathrm{direct}}(t)
    =
    1,
    \qquad
    \forall (i,j)
    \in
    E(R).
\end{equation}

The route-level availability indicator is

\begin{equation}
    \chi_{\mathrm{avail}}(R,t)
    =
    \prod_{
    (i,j)
    \in
    E(R)
    }
    u_{ij}^{\mathrm{direct}}(t).
\end{equation}

If

\begin{equation}
    \chi_{\mathrm{avail}}(R,t)
    =
    0,
\end{equation}

the route is removed from the directly feasible set.

\subsection{Predicted Consumption-Time Calculation}

Therefore, the temporal feasibility evaluation is performed individually for each e-link based on its predicted consumption time. Therefore, the controller estimates the consumption time of each e-link according to the expected swapping schedule.

Let

\begin{equation}
    \widehat{\tau}_{ij}^{\mathrm{cons}}(R,t)
\end{equation}

denote the estimated additional time before e-link $(i,j)$ is consumed as part of route $R$.

A general estimate is

\begin{equation}
    \widehat{\tau}_{ij}^{\mathrm{cons}}(R,t)
    =
    \widehat{D}_{\mathrm{queue}}(R,t)
    +
    \widehat{D}_{\mathrm{coord}}(R,t)
    +
    \widehat{D}_{ij}^{\mathrm{sched}}(R,t),
\end{equation}

where

\begin{itemize}
    \item $\widehat{D}_{\mathrm{queue}}(R,t)$ is the expected request waiting delay;
    \item $\widehat{D}_{\mathrm{coord}}(R,t)$ is the classical coordination delay; and
    \item $\widehat{D}_{ij}^{\mathrm{sched}}(R,t)$ is the delay introduced by the selected swapping schedule.
\end{itemize}

For a simple sequential swapping schedule, the estimated scheduling delay can be approximated as

\begin{equation}
    \widehat{D}_{ij}^{\mathrm{sched}}(R,t)
    =
    \ell_{ij}(R)
    D_{\mathrm{swap}},
\end{equation}

where $\ell_{ij}(R)$ is the number of swapping stages completed before e-link $(i,j)$ is consumed.

For a parallel or nested swapping schedule, this value may be smaller and is determined by the swapping depth of the route.

\subsection{Link-Level Fidelity Test}

The predicted fidelity of e-link $(i,j)$ at its estimated consumption time is

\begin{equation}
    \widehat{F}_{ij}^{\mathrm{cons}}(R,t)
    =
    f_{\mathrm{dec}}
    \left(
    F_{ij}^{\mathrm{eff}}(t),
    \widehat{\tau}_{ij}^{\mathrm{cons}}(R,t),
    T_{2,ij}
    \right).
\end{equation}

The route passes the link-level fidelity test only if

\begin{equation}
    \widehat{F}_{ij}^{\mathrm{cons}}(R,t)
    \geq
    F_{\min}^{\mathrm{link}},
    \qquad
    \forall (i,j)
    \in
    E(R).
\end{equation}

The corresponding indicator is

\begin{equation}
    \chi_{\mathrm{linkF}}(R,t)
    =
    \prod_{
    (i,j)
    \in
    E(R)
    }
    \mathbf{1}
    \left[
    \widehat{F}_{ij}^{\mathrm{cons}}(R,t)
    \geq
    F_{\min}^{\mathrm{link}}
    \right].
\end{equation}

\subsection{End-to-End Fidelity Test}

Using the predicted consumption-time fidelities, the controller estimates the end-to-end route fidelity.

Under the Werner-state approximation, the predicted link parameter is

\begin{equation}
    \widehat{p}_{ij}^{W}(R,t)
    =
    \frac{
    4\widehat{F}_{ij}^{\mathrm{cons}}(R,t)-1
    }{3}.
\end{equation}

The predicted end-to-end Werner parameter is

\begin{equation}
    \widehat{p}_{s,d}^{W,\mathrm{end}}(R,t)
    =
    \left(
    \prod_{
    (i,j)
    \in
    E(R)
    }
    \widehat{p}_{ij}^{W}(R,t)
    \right)
    \left(
    \prod_{
    v_r
    \in
    \mathrm{Int}(R)
    }
    \lambda_r
    \right),
\end{equation}

where $\mathrm{Int}(R)$ denotes the set of intermediate ESPs participating in entanglement swapping along route $R$.

The predicted end-to-end fidelity is

\begin{equation}
    \widehat{F}_{s,d}^{\mathrm{end}}(R,t)
    =
    \frac{
    1
    +
    3\widehat{p}_{s,d}^{W,\mathrm{end}}(R,t)
    }{4}.
\end{equation}

The route passes the end-to-end fidelity test when

\begin{equation}
    \widehat{F}_{s,d}^{\mathrm{end}}(R,t)
    \geq
    F_{\min}^{\mathrm{end}}.
\end{equation}

The associated indicator is

\begin{equation}
    \chi_{\mathrm{endF}}(R,t)
    =
    \mathbf{1}
    \left[
    \widehat{F}_{s,d}^{\mathrm{end}}(R,t)
    \geq
    F_{\min}^{\mathrm{end}}
    \right].
\end{equation}


\subsection{Route-Delay Test}

The predicted route-completion delay is denoted by

\begin{equation}
    \widehat{D}_{s,d}^{\mathrm{route}}(R,t).
\end{equation}

A general decomposition is

\begin{equation}
    \widehat{D}_{s,d}^{\mathrm{route}}(R,t)
    =
    \widehat{D}_{\mathrm{queue}}(R,t)
    +
    \widehat{D}_{\mathrm{coord}}(R,t)
    +
    \widehat{D}_{\mathrm{pur}}(R,t)
    +
    \widehat{D}_{\mathrm{swap}}(R,t).
\end{equation}

The delay estimation is route-dependent because different candidate routes may contain different numbers of swapping operations, purification requirements, and queueing conditions.

The route satisfies the delay constraint only if

\begin{equation}
    \widehat{D}_{s,d}^{\mathrm{route}}(R,t)
    \leq
    D_{\max}.
\end{equation}

The route-delay indicator is

\begin{equation}
    \chi_{\mathrm{delay}}(R,t)
    =
    \mathbf{1}
    \left[
    \widehat{D}_{s,d}^{\mathrm{route}}(R,t)
    \leq
    D_{\max}
    \right].
\end{equation}

\subsection{Memory and E-Bit Availability Test}

The number of stored e-bits on each artificial e-link must be sufficient to satisfy the entanglement request.

For a single-pair request, the link-level resource condition is

\begin{equation}
    n_{ij}^{\mathrm{ebit}}(t)
    \geq
    1,
    \qquad
    \forall (i,j)
    \in
    E(R).
\end{equation}

For a request requiring $b_{s,d}$ entangled pairs, where $b_{s,d}$ denotes the required number of entangled pairs between source and destination, the condition becomes

\begin{equation}
    n_{ij}^{\mathrm{ebit}}(t)
    \geq
    b_{s,d},
    \qquad
    \forall (i,j)
    \in
    E(R).
\end{equation}

The route-level resource indicator is

\begin{equation}
    \chi_{\mathrm{mem}}(R,t)
    =
    \prod_{
    (i,j)
    \in
    E(R)
    }
    \mathbf{1}
    \left[
    n_{ij}^{\mathrm{ebit}}(t)
    \geq
    b_{s,d}
    \right].
\end{equation}


\subsection{Overall Temporal-Feasibility Decision}

The overall temporal-feasibility indicator is defined as

\begin{equation}
    \chi_{s,d}(R,t)
    =
    \chi_{\mathrm{avail}}(R,t)
    \chi_{\mathrm{linkF}}(R,t)
    \chi_{\mathrm{coh}}(R,t)
    \chi_{\mathrm{endF}}(R,t)
    \chi_{\mathrm{delay}}(R,t)
    \chi_{\mathrm{mem}}(R,t).
\end{equation}

A candidate route is directly feasible when

\begin{equation}
    \chi_{s,d}(R,t)
    =
    1.
\end{equation}

The temporally feasible route set is therefore

\begin{equation}
    \mathcal{R}_{s,d}^{\mathrm{TF}}(t)
    =
    \left\{
    R
    \in
    \mathcal{R}_{s,d}^{\mathrm{cand}}(t)
    :
    \chi_{s,d}(R,t)
    =
    1
    \right\}.
\end{equation}

Candidate routes that fail the direct-feasibility test are not immediately discarded if their infeasibility can be corrected through limited regeneration or repair operations.

Such routes are transferred to the route-repair set

\begin{equation}
    \mathcal{R}_{s,d}^{\mathrm{repair}}(t).
\end{equation}

A route enters this set only if its predicted repair delay and repair cost satisfy

\begin{equation}
    \widehat{D}_{s,d}^{\mathrm{repair}}(R,t)
    \leq
    D_{\max}^{\mathrm{repair}}
\end{equation}

and

\begin{equation}
    C_{s,d}^{\mathrm{repair}}(R,t)
    \leq
    C_{\max}^{\mathrm{repair}}.
\end{equation}

This two-stage filtering process distinguishes among immediately feasible, repairable, and infeasible routes. It prevents the route-selection module from assigning low cost to a path that is topologically available but temporally incapable of satisfying the entanglement request.
\section{Time-Dependent Route Selection}

The time-dependent route-selection module chooses the best route from the set of temporally feasible candidates.

For an entanglement request between source ESP $v_s$ and destination ESP $v_d$, the feasible candidate set is

\begin{equation}
    \mathcal{R}_{s,d}^{\mathrm{TF}}(t).
\end{equation}

For each candidate route

\begin{equation}
    R
    \in
    \mathcal{R}_{s,d}^{\mathrm{TF}}(t),
\end{equation}

the controller evaluates the complete time-dependent route cost.

The route cost is defined as

\begin{equation}
    w_t(R)
    =
    \sum_{
    (i,j)
    \in
    E(R)
    }
    w_t(i,j)
    +
    C_{\mathrm{route}}^{\mathrm{swap}}(R,t)
    +
    C_{\mathrm{route}}^{\mathrm{delay}}(R,t),
\end{equation}

where

\begin{itemize}
    \item $w_t(i,j)$ is the time-dependent cost of artificial e-link $(i,j)$;
    \item $C_{\mathrm{route}}^{\mathrm{swap}}(R,t)$ is the route-level swapping cost; and
    \item $C_{\mathrm{route}}^{\mathrm{delay}}(R,t)$ is the normalized completion-delay cost.
\end{itemize}

The link-level time-dependent cost is

\begin{equation}
    w_t(i,j)
    =
    \alpha_F C_F(i,j,R,t)
    +
    \alpha_A C_A(i,j,t)
    +
    \alpha_B C_B(i,j,R,t)
    +
    \alpha_R C_R(i,j,t)
    +
    \alpha_S C_S(i,j).
\end{equation}

The coefficients satisfy

\begin{equation}
    \alpha_F
    +
    \alpha_A
    +
    \alpha_B
    +
    \alpha_R
    +
    \alpha_S
    =
    1.
\end{equation}

The fidelity cost is defined as

\begin{equation}
    C_F(i,j,R,t)
    =
    1
    -
    \widehat{F}_{ij}^{\mathrm{cons}}(R,t).
\end{equation}

This term penalizes e-links that are expected to have low fidelity when they are consumed.

The age cost is

\begin{equation}
    C_A(i,j,t)
    =
    \min
    \left\{
    \frac{
    A_{ij}(t)
    }{
    A_{ij}^{\max}
    },
    1
    \right\}.
\end{equation}

The normalized coherence-budget cost is defined as

\begin{equation}
    C_B(i,j,R,t)
    =
    \frac{
    1
    }{
    \widetilde{B}_{ij}(R,t)
    +
    \varepsilon
    },
\end{equation}

where

\begin{equation}
    \widetilde{B}_{ij}(R,t)
    =
    \frac{
    B_{ij}(R,t)
    }{
    T_{2,ij}
    }.
\end{equation}

Since only temporally feasible routes reach this module, links with non-positive coherence budget are already removed before route selection:

\begin{equation}
    \widetilde{B}_{ij}(R,t)
    >
    0.
\end{equation}

The normalized regeneration-risk cost is

\begin{equation}
    C_R(i,j,t)
    =
    \lambda_D
    \frac{
    D_{ij}^{\mathrm{regen}}
    }{
    D_{\max}^{\mathrm{regen}}
    }
    +
    \lambda_C
    \frac{
    C_{ij}^{\mathrm{regen}}
    }{
    C_{\max}^{\mathrm{regen}}
    },
\end{equation}

where

\begin{equation}
    \lambda_D+\lambda_C=1.
\end{equation}

The structural routing cost is represented by

\begin{equation}
    C_S(i,j).
\end{equation}

This component distinguishes local e-neighbor links, anchor links, and tracked long-range links according to their resource and maintenance requirements.

\subsection{Route-Level Cost Components}

The swapping cost of route $R$ is defined as

\begin{equation}
    C_{\mathrm{route}}^{\mathrm{swap}}(R,t)
    =
    \alpha_Q
    \sum_{
    v_r
    \in
    \mathrm{Int}(R)
    }
    \left(
    1-q_r
    \right),
\end{equation}

where $q_r$ is the swapping-success probability at intermediate ESP $v_r$.

The normalized route-delay cost is

\begin{equation}
    C_{\mathrm{route}}^{\mathrm{delay}}(R,t)
    =
    \alpha_D
    \frac{
    \widehat{D}_{s,d}^{\mathrm{route}}(R,t)
    }{
    D_{\max}
    }.
\end{equation}

The selected directly feasible route is

\begin{equation}
    R_{s,d}^{*}(t)
    =
    \arg\min_{
    R
    \in
    \mathcal{R}_{s,d}^{\mathrm{TF}}(t)
    }
    w_t(R).
\end{equation}

When multiple routes have the same minimum cost, the controller applies the following tie-breaking order:

\begin{enumerate}
    \item select the route with the highest predicted end-to-end fidelity;
    \item select the route with the largest minimum coherence budget;
    \item select the route with the smallest predicted completion delay;
    \item select the route with the fewest artificial e-links.
\end{enumerate}

The minimum route-level coherence budget is

\begin{equation}
    B_{\min}(R,t)
    =
    \min_{
    (i,j)
    \in
    E(R)
    }
    B_{ij}(R,t).
\end{equation}

The route-selection decision can therefore be represented lexicographically as

\begin{equation}
    \min
    \left[
    w_t(R),
    -
    \widehat{F}_{s,d}^{\mathrm{end}}(R,t),
    -
    B_{\min}(R,t),
    \widehat{D}_{s,d}^{\mathrm{route}}(R,t),
    |E(R)|
    \right].
\end{equation}

After selecting a route, the required e-bits are temporarily reserved before swapping begins.

For a request requiring $b_{s,d}$ entangled pairs, the reservation variable is

\begin{equation}
    z_{ij}^{s,d}(t)
    =
    \begin{cases}
        b_{s,d},
        & \text{if $(i,j)$ belongs to the selected route},\\
        0,
        & \text{otherwise}.
    \end{cases}
\end{equation}

The reservation constraint is

\begin{equation}
    \sum_{
    (s,d)
    \in
    \mathcal{D}(t)
    }
    z_{ij}^{s,d}(t)
    \leq
    n_{ij}^{\mathrm{ebit}}(t).
\end{equation}

This reservation prevents the same stored e-bit from being assigned to multiple requests.

After successful end-to-end swapping, the consumed e-bits are removed:

\begin{equation}
    n_{ij}^{\mathrm{ebit}}(t+1)
    =
    n_{ij}^{\mathrm{ebit}}(t)
    -
    z_{ij}^{s,d}(t).
\end{equation}

If a swapping operation fails, all e-bits already consumed by the failed route are treated as lost.

The route outcome is represented by

\begin{equation}
    Y_{s,d}(t)
    =
    \begin{cases}
        1, & \text{if end-to-end entanglement is established},\\
        0, & \text{otherwise}.
    \end{cases}
\end{equation}

The realized route success probability is estimated as

\begin{equation}
    P_{\mathrm{succ}}(R,t)
    =
    \left(
    \prod_{
    (i,j)
    \in
    E(R)
    }
    p_{ij}^{\mathrm{avail}}(t)
    \right)
    \left(
    \prod_{
    v_r
    \in
    \mathrm{Int}(R)
    }
    q_r
    \right),
\end{equation}

where $p_{ij}^{\mathrm{avail}}(t)$ represents the probability that artificial e-link $(i,j)$ remains usable until its scheduled consumption time.

A reliability-aware extension of the route cost is

\begin{equation}
    w_t^{\mathrm{rel}}(R)
    =
    w_t(R)
    +
    \alpha_P
    \left[
    -\log
    P_{\mathrm{succ}}(R,t)
    \right].
\end{equation}

This logarithmic transformation converts the multiplicative success probability into an additive routing penalty.

The route selected under the reliability-aware metric is

\begin{equation}
    R_{s,d}^{*,\mathrm{rel}}(t)
    =
    \arg\min_{
    R
    \in
    \mathcal{R}_{s,d}^{\mathrm{TF}}(t)
    }
    w_t^{\mathrm{rel}}(R).
\end{equation}

The reliability-aware metric provides an optional extension of the basic temporal entangling-cost metric. Its effectiveness will be evaluated through comparative experiments.

When

\begin{equation}
    \mathcal{R}_{s,d}^{\mathrm{TF}}(t)
    =
    \emptyset,
\end{equation}

the module does not select a directly feasible route.

Instead, it forwards the request and the repairable candidate set

\begin{equation}
    \mathcal{R}_{s,d}^{\mathrm{repair}}(t)
\end{equation}

to the proactive regeneration and route-repair process described in the following section.

The normalized fidelity risk is defined as

\begin{equation}
    \widehat{C}_{F}(i,j,t)
    =
    \min
    \left\{
    \frac{
    F_{\mathrm{high}}
    -
    \widehat{F}_{ij}^{\mathrm{cons}}(t)
    }{
    F_{\mathrm{high}}
    -
    F_{\min}^{\mathrm{link}}
    +
    \varepsilon
    },
    1
    \right\}.
\end{equation}

The predicted demand relevance may be estimated from the number of active or expected routes that can use the e-link:

\begin{equation}
    Q_{ij}^{\mathrm{demand}}(t)
    =
    \frac{
    N_{ij}^{\mathrm{demand}}(t)
    }{
    \max_{(u,v)}
    N_{uv}^{\mathrm{demand}}(t)
    +
    \varepsilon
    }.
\end{equation}

The normalized regeneration cost is

\begin{equation}
    \widehat{C}_{R}(i,j,t)
    =
    \lambda_D
    \frac{
    D_{ij}^{\mathrm{regen}}
    }{
    D_{\max}^{\mathrm{regen}}
    }
    +
    \lambda_C
    \frac{
    C_{ij}^{\mathrm{regen}}
    }{
    C_{\max}^{\mathrm{regen}}
    },
\end{equation}

where

\begin{equation}
    \lambda_D+\lambda_C=1.
\end{equation}

A link becomes a proactive-regeneration candidate when

\begin{equation}
    \Pi_{ij}^{\mathrm{regen}}(t)
    \geq
    \Pi_{\min}
\end{equation}

and sufficient regeneration resources are available.

The candidate set is

\begin{equation}
    \mathcal{E}_{\mathrm{regen}}^{\mathrm{cand}}(t)
    =
    \left\{
    (i,j)
    \in
    \mathcal{E}_{\mathrm{crit}}(t)
    :
    \Pi_{ij}^{\mathrm{regen}}(t)
    \geq
    \Pi_{\min}
    \right\}.
\end{equation}

The regeneration-decision variable is

\begin{equation}
    r_{ij}(t)
    \in
    \{0,1\}.
\end{equation}

It is defined as

\begin{equation}
    r_{ij}(t)
    =
    \begin{cases}
        1, & \text{if e-link $(i,j)$ is selected for regeneration},\\
        0, & \text{otherwise}.
    \end{cases}
\end{equation}

The controller selects at most $R_{\max}$ highest-priority regeneration candidates:

\begin{equation}
    \mathcal{E}_{\mathrm{regen}}^{*}(t)
    =
    \operatorname{TopK}
    \left(
    \mathcal{E}_{\mathrm{regen}}^{\mathrm{cand}}(t),
    R_{\max}
    \right).
\end{equation}

The selected regeneration actions must satisfy the regeneration budget constraint:

\begin{equation}
    |
    \mathcal{E}_{\mathrm{regen}}^{*}(t)
    |
    \leq
    R_{\max}.
\end{equation}

Let $m_i(t)$ denote the currently occupied quantum-memory capacity at ESP $v_i$. The regeneration decision must satisfy

\begin{equation}
    m_i(t)
    +
    \sum_{
    j:
    (i,j)
    \in
    \mathcal{E}_{\mathrm{regen}}^{*}(t)
    }
    r_{ij}(t)
    \leq
    M_i,
    \qquad
    \forall v_i
    \in
    V_2.
\end{equation}

TThey must also satisfy the physical-generation capacity constraint:

\begin{equation}
    \sum_{
    (i,j)
    \in
    E_{\mathrm{A}}(t)
    }
    r_{ij}(t)
    g_{ij}^{\mathrm{req}}
    \leq
    G_{\max},
\end{equation}

where $g_{ij}^{\mathrm{req}}$ represents the expected number of elementary-link generation attempts required for regenerating artificial e-link $(i,j)$, and $G_{\max}$ denotes the available generation-attempt budget.

For an artificial e-link generated through physical path $P_{ij}^{\mathrm{gen}}$, the expected regeneration delay is

\begin{equation}
    D_{ij}^{\mathrm{regen}}
    =
    D^{\mathrm{phys}}
    \left(
    P_{ij}^{\mathrm{gen}}
    \right)
    +
    D_{ij}^{\mathrm{swap}}
    +
    D_{ij}^{\mathrm{coord}}.
\end{equation}

When purification is required, the regeneration delay becomes

\begin{equation}
    D_{ij}^{\mathrm{regen,pur}}
    =
    D_{ij}^{\mathrm{regen}}
    +
    D_{ij}^{\mathrm{pur}}.
\end{equation}

The regeneration process may follow either a replace-after-success strategy or a discard-before-generation strategy.

Under the replace-after-success strategy, the old artificial e-link remains available until the new e-link is successfully generated. After successful regeneration,

\begin{equation}
    F_{ij}^{\mathrm{eff}}(t+1)
    =
    F_{ij}^{\mathrm{A},0},
\end{equation}

\begin{equation}
    A_{ij}(t+1)
    =
    0,
\end{equation}

and

\begin{equation}
    B_{ij}(t+1)
    =
    T_{2,ij}
    -
    \widehat{D}_{ij}^{\mathrm{use}}(t+1).
\end{equation}

This strategy reduces temporary disconnection risk but requires additional memory capacity because both old and newly generated e-links may coexist during the transition period.

Under the discard-before-generation strategy, the old e-link is released before regeneration. This reduces memory consumption but temporarily removes the link from the usable overlay.

The replacement-policy variable is defined as

\begin{equation}
    \rho_{ij}^{\mathrm{replace}}
    \in
    \{0,1\},
\end{equation}

where

\begin{equation}
    \rho_{ij}^{\mathrm{replace}}
    =
    1
\end{equation}

indicates that the replace-after-success strategy is selected.

The policy is determined according to available memory capacity:

\begin{equation}
    \rho_{ij}^{\mathrm{replace}}
    =
    \mathbf{1}
    \left[
    M_i-m_i(t)
    \geq
    1
    \ \land\
    M_j-m_j(t)
    \geq
    1
    \right].
\end{equation}

Regeneration is considered successful only when the required elementary-link generation, entanglement swapping, and coordination operations are completed successfully.

Let

\begin{equation}
    Z_{ij}^{\mathrm{regen}}(t)
    \in
    \{0,1\}
\end{equation}

denote the regeneration outcome of artificial e-link $(i,j)$ at time slot $t$.

The regeneration success probability is approximated as

\begin{equation}
    P
    \left[
    Z_{ij}^{\mathrm{regen}}(t)=1
    \right]
    =
    P_{ij}^{\mathrm{regen}}
    =
    P_{\mathrm{succ}}
    \left(
    P_{ij}^{\mathrm{gen}}
    \right).
\end{equation}

If regeneration fails, the e-link remains in its previous state under the replace-after-success strategy. Under the discard-before-generation strategy, the e-link remains unavailable until a subsequent regeneration attempt succeeds.

Reactive route repair uses the same regeneration mechanism but assigns additional priority according to the blocked request.

For a repairable route $R$, the repair priority of a critical e-link is increased by

\begin{equation}
    \Pi_{ij}^{\mathrm{repair}}(R,t)
    =
    \Pi_{ij}^{\mathrm{regen}}(t)
    +
    \beta_{\mathrm{req}}
    I_{ij}^{R},
\end{equation}

where

\begin{equation}
    I_{ij}^{R}
    =
    \mathbf{1}
    \left[
    (i,j)
    \in
    \mathcal{E}_{\mathrm{crit}}(R)
    \right].
\end{equation}

This request-specific term provides higher priority to links whose regeneration can restore a currently infeasible route.

A regeneration action is approved only when its expected benefit exceeds its expected cost.

The expected regeneration utility is defined as

\begin{equation}
    U_{ij}^{\mathrm{regen}}(t)
    =
    P_{ij}^{\mathrm{regen}}
    G_{ij}^{\mathrm{future}}(t)
    -
    \kappa_D
    D_{ij}^{\mathrm{regen}}
    -
    \kappa_C
    C_{ij}^{\mathrm{regen}},
\end{equation}

where $G_{ij}^{\mathrm{future}}(t)$ represents the predicted future routing benefit obtained by refreshing artificial e-link $(i,j)$.

Regeneration is executed only if

\begin{equation}
    U_{ij}^{\mathrm{regen}}(t)
    >
    0.
\end{equation}

This condition prevents unnecessary regeneration of low-impact e-links whose expected routing benefit does not justify the required generation resources and delay.

The proposed proactive-regeneration mechanism therefore balances temporal freshness, fidelity preservation, structural importance, future demand, memory availability, physical-generation capacity, and regeneration overhead.
\section{Overlay Maintenance Strategy}

The overlay-maintenance strategy updates the artificial entanglement overlay according to temporal degradation, routing demand, memory availability, and compact-routing constraints.

Its purpose is not to keep every artificial e-link permanently fresh. Instead, the strategy maintains a bounded set of useful e-links whose expected routing benefit justifies their generation and storage cost.

At time slot $t$, the stored artificial overlay is

\begin{equation}
    G_{\mathrm{A}}(t)
    =
    \left(
    V_2,
    E_{\mathrm{A}}(t)
    \right).
\end{equation}

The overlay is updated through four principal actions:

\begin{equation}
    \mathrm{Retain},
    \qquad
    \mathrm{Regenerate},
    \qquad
    \mathrm{Replace},
    \qquad
    \mathrm{Discard}.
\end{equation}

For each artificial e-link $(i,j)$, the maintenance action is denoted by

\begin{equation}
    a_{ij}^{\mathrm{maint}}(t)
    \in
    \left\{
    \mathrm{retain},
    \mathrm{regenerate},
    \mathrm{replace},
    \mathrm{discard}
    \right\}.
\end{equation}


\subsection{Retention Decision}

An e-link is retained without regeneration when it remains temporally reliable and sufficiently useful.

The retention indicator is

\begin{equation}
    k_{ij}^{\mathrm{retain}}(t)
    =
    \mathbf{1}
    \left[
    S_{ij}^{\mathrm{state}}(t)
    \in
    \{
    \mathrm{Safe},
    \mathrm{Warning}
    \}
    \right].
\end{equation}

A warning-state link may be retained only if its predicted consumption fidelity and remaining coherence budget satisfy the minimum operational margins:

\begin{equation}
    \widehat{F}_{ij}^{\mathrm{cons}}(t)
    \geq
    F_{\min}^{\mathrm{link}}
\end{equation}

and

\begin{equation}
    B_{ij}(t)
    \geq
    \delta_{\mathrm{safe}}.
\end{equation}


\subsection{Regeneration Decision}

A critical e-link is considered for regeneration when its expected future routing benefit is high and its regeneration cost is acceptable.

The regeneration decision follows the previously defined variable:

\begin{equation}
    r_{ij}(t)
    =
    \mathbf{1}
    \left[
    \Pi_{ij}^{\mathrm{regen}}(t)
    \geq
    \Pi_{\min}
    \right],
\end{equation}

subject to memory, physical-generation, and per-slot regeneration constraints.

The total number of regeneration actions satisfies

\begin{equation}
    \sum_{(i,j)\in E_{\mathrm{A}}(t)}
    r_{ij}(t)
    \leq
    R_{\max}.
\end{equation}


\subsection{Replacement Decision}

Replacement changes the endpoint pair or physical realization of an artificial e-link when repeated regeneration becomes inefficient or when the existing link has limited routing value.

Let

\begin{equation}
    h_{ij}^{\mathrm{fail}}(t)
\end{equation}

denote the recent number of failed regeneration attempts for e-link $(i,j)$.

The link becomes a replacement candidate when

\begin{equation}
    h_{ij}^{\mathrm{fail}}(t)
    \geq
    H_{\max}^{\mathrm{fail}},
\end{equation}

or when

\begin{equation}
    U_{ij}^{\mathrm{regen}}(t)
    \leq
    0.
\end{equation}

A replacement candidate may also be identified when its structural importance is low while another non-stored e-link provides higher predicted benefit.

Let

\begin{equation}
    \mathcal{E}_{\mathrm{new}}^{\mathrm{cand}}(t)
\end{equation}

denote the set of candidate artificial e-links that are not currently stored.

For a candidate new e-link $(u,v)$, the predicted placement utility is defined as

\begin{equation}
    U_{uv}^{\mathrm{place}}(t)
    =
    \theta_D
    Q_{uv}^{\mathrm{demand}}(t)
    +
    \theta_S
    I_{uv}^{\mathrm{struct}}
    +
    \theta_C
    G_{uv}^{\mathrm{connect}}(t)
    -
    \theta_R
    \widehat{C}_{R}(u,v,t).
\end{equation}

The weights satisfy

\begin{equation}
    \theta_D
    +
    \theta_S
    +
    \theta_C
    +
    \theta_R
    =
    1.
\end{equation}

A stored e-link $(i,j)$ may be replaced by candidate $(u,v)$ when

\begin{equation}
    U_{uv}^{\mathrm{place}}(t)
    -
    U_{ij}^{\mathrm{retain}}(t)
    \geq
    \Delta U_{\min},
\end{equation}

where $U_{ij}^{\mathrm{retain}}(t)$ represents the predicted utility of keeping the existing artificial e-link without replacement.

The threshold $\Delta U_{\min}$ prevents frequent overlay changes caused by small utility fluctuations.


\subsection{Discard Decision}

An artificial e-link is discarded when it becomes expired, repeatedly unused, structurally redundant, or too expensive to maintain.

The discard indicator is

\begin{equation}
    d_{ij}(t)
    \in
    \{0,1\}.
\end{equation}

It is defined as

\begin{equation}
    d_{ij}(t)
    =
    \mathbf{1}
    \left[
    S_{ij}^{\mathrm{state}}(t)
    =
    \mathrm{Expired}
    \right].
\end{equation}

A non-expired link may also be discarded when the following redundancy conditions are simultaneously satisfied:

\begin{equation}
    H_{ij}^{\mathrm{use}}(t)
    <
    H_{\min}^{\mathrm{use}},
\end{equation}

\begin{equation}
    I_{ij}^{\mathrm{struct}}
    <
    I_{\min}^{\mathrm{struct}},
\end{equation}

and

\begin{equation}
    U_{ij}^{\mathrm{regen}}(t)
    \leq
    0.
\end{equation}

An e-link is not discarded when doing so would violate a required compact-overlay connectivity condition.

\subsection{Compact-Overlay Constraints}

For every ESP $v_i$, the number of maintained routing-table entries is bounded by

\begin{equation}
    |\mathcal{T}_i(t)|
    \leq
    T_i^{\max},
\end{equation}

where

\begin{equation}
    T_i^{\max}
    =
    c_{\mathrm{tab}}
    \sqrt{n_e}
    \log n_e .
\end{equation}

This constraint preserves the compact-routing property by preventing each ESP from maintaining complete global routing information.

In the Partial-Anchor Scheme, the required local and anchor connectivity structures must be preserved:

\begin{equation}
    E_{\mathrm{local}}^{\mathrm{req}}
    \subseteq
    E_{\mathrm{A}}^{\mathrm{PA}}(t).
\end{equation}

In addition, the anchor overlay subgraph should remain connected whenever sufficient physical entanglement resources are available.

In the Full-Anchor Scheme, each ESP maintains a bounded tracked-link set:

\begin{equation}
    \left|
    \mathcal{L}_{i}^{\mathrm{track}}(t)
    \right|
    \leq
    L_{\max}^{\mathrm{track}}.
\end{equation}

New tracked links are admitted only when available routing-table capacity exists or when an existing tracked link is removed.


\subsection{Memory-Aware Maintenance}

The overlay-maintenance strategy must satisfy the quantum-memory constraint

\begin{equation}
    \sum_{
    j:
    (i,j)
    \in
    E_{\mathrm{A}}(t)
    }
    n_{ij}^{\mathrm{ebit}}(t)
    \leq
    M_i,
    \qquad
    \forall v_i
    \in
    V_2 .
\end{equation}

When memory capacity becomes insufficient, the controller releases artificial e-links according to a discard-priority score.

The discard-priority score is defined as

\begin{equation}
    \Pi_{ij}^{\mathrm{drop}}(t)
    =
    \phi_1
    \left(
    1-H_{ij}^{\mathrm{use}}(t)
    \right)
    +
    \phi_2
    \left(
    1-I_{ij}^{\mathrm{struct}}
    \right)
    +
    \phi_3
    \widehat{C}_{R}(i,j,t)
    +
    \phi_4
    I_{ij}^{\mathrm{expired}}(t),
\end{equation}

where

\begin{equation}
    I_{ij}^{\mathrm{expired}}(t)
    =
    \mathbf{1}
    \left[
    S_{ij}^{\mathrm{state}}(t)
    =
    \mathrm{Expired}
    \right].
\end{equation}

Links with higher discard priority are released first while preserving required compact-overlay connectivity.


\subsection{Maintenance Triggering}

Overlay maintenance may be executed periodically or through event-driven triggers.

Periodic maintenance is performed every

\begin{equation}
    \tau_{\mathrm{maint}}
\end{equation}

time slots.

Event-driven maintenance is triggered when one of the following conditions occurs:

\begin{equation}
    S_{ij}^{\mathrm{state}}(t)
    =
    \mathrm{Critical}
\end{equation}

for a structurally important e-link,

\begin{equation}
    S_{ij}^{\mathrm{state}}(t)
    =
    \mathrm{Expired},
\end{equation}

\begin{equation}
    m_i(t)
    \geq
    M_i,
\end{equation}

or

\begin{equation}
    P_{\mathrm{disc}}(t)
    \geq
    P_{\mathrm{disc}}^{\max}.
\end{equation}

To prevent excessive overlay reconfiguration, a minimum holding interval

\begin{equation}
    \tau_{\mathrm{hold}}
\end{equation}

is enforced between two structural changes affecting the same routing-table entry.

\subsection{Overlay-State Update}

Let

\begin{equation}
    E_{\mathrm{retain}}(t),
    \quad
    E_{\mathrm{regen}}^{\mathrm{succ}}(t),
    \quad
    E_{\mathrm{new}}^{\mathrm{succ}}(t),
    \quad
    E_{\mathrm{drop}}(t)
\end{equation}

denote the sets of retained, successfully regenerated, newly established, and discarded artificial e-links, respectively.

The updated artificial overlay edge set is defined as

\begin{equation}
    E_{\mathrm{A}}(t+1)
    =
    \left(
    E_{\mathrm{retain}}(t)
    \cup
    E_{\mathrm{regen}}^{\mathrm{succ}}(t)
    \cup
    E_{\mathrm{new}}^{\mathrm{succ}}(t)
    \right)
    \setminus
    E_{\mathrm{drop}}(t).
\end{equation}

The corresponding routing tables are updated only after successful entanglement generation or regeneration is confirmed through classical heralding information.

Therefore, a routing entry is not advertised as usable until the corresponding artificial e-link has been successfully established and its initial temporal state has been recorded.

The proposed maintenance strategy balances temporal reliability and resource efficiency by preserving useful links, selectively regenerating critical links, replacing inefficient links, and removing stale or low-value links while maintaining the compact structure of the artificial overlay.


\section{Proposed Routing Algorithm}

The proposed routing algorithm integrates temporal-state monitoring, compact candidate construction, temporal-feasibility filtering, time-dependent route selection, and selective artificial e-link regeneration.

The algorithm is executed at each time slot $t$ and processes the active entanglement-request set

\begin{equation}
    \mathcal{D}(t).
\end{equation}

The main inputs are

\begin{equation}
    G_{\mathrm{P}},
    \quad
    G_{\mathrm{A}}(t),
    \quad
    \mathcal{D}(t),
    \quad
    \mathcal{T}(t),
    \quad
    \mathcal{M}(t),
\end{equation}

together with fidelity, coherence, delay, memory, and regeneration parameters.

The main outputs are

\begin{equation}
    \mathcal{Y}(t),
    \quad
    \mathcal{E}_{\mathrm{regen}}^{*}(t),
    \quad
    \mathcal{E}_{\mathrm{drop}}(t),
    \quad
    G_{\mathrm{A}}(t+1),
\end{equation}

where $\mathcal{Y}(t)$ contains selected routes and their execution outcomes.

The complete procedure is described as follows.


\subsection{Algorithmic Procedure}

\begin{enumerate}

\item \textbf{Update temporal e-link states.}

For every artificial e-link $(i,j)\in E_{\mathrm{A}}(t)$, update

\begin{equation}
    A_{ij}(t),
    \quad
    F_{ij}^{\mathrm{eff}}(t),
    \quad
    \widehat{F}_{ij}^{\mathrm{cons}}(t),
    \quad
    B_{ij}(t),
    \quad
    H_{ij}^{\mathrm{use}}(t).
\end{equation}

Each e-link is classified as safe, warning, critical, or expired according to the temporal-state classification rules.


\item \textbf{Construct the temporally usable overlay.}

The usable edge set is constructed as

\begin{equation}
    E_{\mathrm{U}}(t)
    =
    \left\{
    (i,j)
    \in
    E_{\mathrm{A}}(t)
    :
    S_{ij}^{\mathrm{state}}(t)
    \in
    \{
    \mathrm{Safe},
    \mathrm{Warning}
    \}
    \right\}.
\end{equation}

The temporally usable overlay is then defined as

\begin{equation}
    G_{\mathrm{U}}(t)
    =
    \left(
    V_2,
    E_{\mathrm{U}}(t)
    \right).
\end{equation}


\item \textbf{Process each entanglement request.}

For every request

\begin{equation}
    (s,d)
    \in
    \mathcal{D}(t),
\end{equation}

the controller generates a compact candidate-route set

\begin{equation}
    \mathcal{R}_{s,d}^{\mathrm{cand}}(t).
\end{equation}

The candidate set contains no more than $K_{\max}$ routes.


\item \textbf{Evaluate direct temporal feasibility.}

For every candidate route $R$, calculate

\begin{equation}
    \chi_{\mathrm{avail}}(R,t),
    \quad
    \chi_{\mathrm{linkF}}(R,t),
    \quad
    \chi_{\mathrm{coh}}(R,t),
\end{equation}

and

\begin{equation}
    \chi_{\mathrm{endF}}(R,t),
    \quad
    \chi_{\mathrm{delay}}(R,t),
    \quad
    \chi_{\mathrm{mem}}(R,t).
\end{equation}

    The overall feasibility indicator is

\begin{equation}
    \chi_{s,d}(R,t)
    =
    \chi_{\mathrm{avail}}(R,t)
    \chi_{\mathrm{linkF}}(R,t)
    \chi_{\mathrm{coh}}(R,t)
    \chi_{\mathrm{endF}}(R,t)
    \chi_{\mathrm{delay}}(R,t)
    \chi_{\mathrm{mem}}(R,t).
\end{equation}


\item \textbf{Build the directly feasible route set.}

The temporally feasible route set is defined as

\begin{equation}
    \mathcal{R}_{s,d}^{\mathrm{TF}}(t)
    =
    \left\{
    R
    \in
    \mathcal{R}_{s,d}^{\mathrm{cand}}(t)
    :
    \chi_{s,d}(R,t)=1
    \right\}.
\end{equation}


\item \textbf{Select a directly feasible route.}

If

\begin{equation}
    \mathcal{R}_{s,d}^{\mathrm{TF}}(t)
    \neq
    \emptyset,
\end{equation}

the controller calculates the complete time-dependent cost

\begin{equation}
    W_t(R)
\end{equation}

for every feasible route.

The selected route is

\begin{equation}
    R_{s,d}^{*}(t)
    =
    \arg\min_{
    R
    \in
    \mathcal{R}_{s,d}^{\mathrm{TF}}(t)
    }
    W_t(R).
\end{equation}

The required e-bits are reserved and the planned swapping operations are executed.


\item \textbf{Construct the repairable route set.}

If no directly feasible route exists, candidate routes containing a bounded number of regenerable critical links are considered.

The repairable route set is defined as

\begin{equation}
    \mathcal{R}_{s,d}^{\mathrm{repair}}(t)
    =
    \left\{
    R
    \in
    \mathcal{R}_{s,d}^{\mathrm{cand}}(t)
    :
    \widehat{D}_{s,d}^{\mathrm{repair}}(R,t)
    \leq
    D_{\max}^{\mathrm{repair}}
    \land
    C_{s,d}^{\mathrm{repair}}(R,t)
    \leq
    C_{\max}^{\mathrm{repair}}
    \right\}.
\end{equation}


\item \textbf{Select a repairable route.}

If

\begin{equation}
    \mathcal{R}_{s,d}^{\mathrm{repair}}(t)
    \neq
    \emptyset,
\end{equation}

the repair-aware cost is calculated as

\begin{equation}
    W_t^{\mathrm{repair}}(R)
    =
    W_t(R)
    +
    \rho_D
    \widehat{D}_{s,d}^{\mathrm{repair}}(R,t)
    +
    \rho_C
    C_{s,d}^{\mathrm{repair}}(R,t).
\end{equation}

The selected repairable route is

\begin{equation}
    R_{s,d}^{\mathrm{repair},*}(t)
    =
    \arg\min_{
    R
    \in
    \mathcal{R}_{s,d}^{\mathrm{repair}}(t)
    }
    W_t^{\mathrm{repair}}(R).
\end{equation}

The critical e-links of this route are added to the request-specific regeneration queue.


\item \textbf{Reject or postpone infeasible requests.}

If both

\begin{equation}
    \mathcal{R}_{s,d}^{\mathrm{TF}}(t)
    =
    \emptyset
\end{equation}

and

\begin{equation}
    \mathcal{R}_{s,d}^{\mathrm{repair}}(t)
    =
    \emptyset,
\end{equation}

the request is classified as temporarily infeasible.

Depending on the request deadline, the controller either postpones or rejects the request.


\item \textbf{Calculate proactive regeneration priorities.}

For every critical artificial e-link, calculate

\begin{equation}
    \Pi_{ij}^{\mathrm{regen}}(t).
\end{equation}

Request-specific critical links receive an additional repair-priority term.


\item \textbf{Select regeneration actions.}

Regeneration candidates are ranked according to priority and expected utility.

The selected regeneration set is

\begin{equation}
    \mathcal{E}_{\mathrm{regen}}^{*}(t),
\end{equation}

subject to

\begin{equation}
    \left|
    \mathcal{E}_{\mathrm{regen}}^{*}(t)
    \right|
    \leq
    R_{\max},
\end{equation}

memory constraints, and physical-generation capacity constraints.


\item \textbf{Execute overlay-maintenance actions.}

Useful links are retained, selected critical links are regenerated, inefficient links are replaced, and expired or low-value links are discarded.


\item \textbf{Update the artificial overlay.}

After receiving generation and swapping outcomes, update

\begin{equation}
    E_{\mathrm{A}}(t+1)
    =
    \left(
    E_{\mathrm{retain}}(t)
    \cup
    E_{\mathrm{regen}}^{\mathrm{succ}}(t)
    \cup
    E_{\mathrm{new}}^{\mathrm{succ}}(t)
    \right)
    \setminus
    E_{\mathrm{drop}}(t).
\end{equation}


\item \textbf{Update compact routing tables.}

A routing-table entry is added only after the corresponding artificial e-link is successfully established and confirmed through classical communication.

Entries associated with consumed, discarded, or permanently unavailable e-links are removed.
\subsection{Algorithm Characteristics}

The proposed algorithm is heuristic rather than globally optimal. It does not solve the complete joint optimization problem involving route selection, regeneration scheduling, and overlay maintenance at every time slot.

Instead, the computational cost is reduced through

\begin{itemize}
    \item compact routing-state restriction;
    \item bounded candidate-route generation;
    \item early temporal-feasibility filtering;
    \item additive time-dependent route-cost evaluation;
    \item bounded regeneration operations; and
    \item threshold-based overlay maintenance.
\end{itemize}

For a fixed network state, parameter configuration, candidate ordering, and tie-breaking rule, the algorithm produces deterministic routing decisions.

Randomness appears only from physical entanglement generation processes, decoherence behavior, and probabilistic swapping operations.

The performance of the proposed method is evaluated by comparison with static compact routing, fidelity-only routing, age-only routing, shortest-path routing, and regeneration-disabled variants.


\section{Computational Complexity Analysis}

The computational complexity of the proposed method consists of temporal-state monitoring, candidate-route construction, temporal-feasibility evaluation, route-cost calculation, and regeneration-priority computation.

Let

\begin{equation}
    n
    =
    |V_2|
\end{equation}

denote the number of Entanglement Service Providers (ESPs), and let

\begin{equation}
    m_{\mathrm{A}}(t)
    =
    |E_{\mathrm{A}}(t)|
\end{equation}

denote the number of stored artificial e-links at time slot $t$.

The number of active entanglement requests at time $t$ is

\begin{equation}
    D_t
    =
    |\mathcal{D}(t)|.
\end{equation}

The maximum number of candidate routes generated for each request is denoted by

\begin{equation}
    K_{\max},
\end{equation}

while the maximum number of artificial e-links contained in one candidate route is denoted by

\begin{equation}
    H_{\max}.
\end{equation}


\subsection{Temporal-State Monitoring Complexity}

The temporal-state monitoring module updates the age, effective fidelity, predicted consumption fidelity, remaining coherence budget, utilization score, and state classification of every stored artificial e-link.

Since each update requires a constant number of operations, the monitoring complexity is

\begin{equation}
    O
    \left(
    m_{\mathrm{A}}(t)
    \right).
\end{equation}

The memory requirement for storing temporal-state information is also

\begin{equation}
    O
    \left(
    m_{\mathrm{A}}(t)
    \right),
\end{equation}

because one temporal-state record is maintained for each artificial e-link.


\subsection{Candidate-Route Construction Complexity}

Candidate-route construction is performed on the temporally usable compact overlay rather than the complete physical network.

When Dijkstra's algorithm or a bounded $K$-shortest-path method is applied to the usable overlay, the shortest-path computation using adjacency lists and a binary heap has complexity

\begin{equation}
    O
    \left(
    m_{\mathrm{U}}(t)
    \log n
    \right),
\end{equation}

where

\begin{equation}
    m_{\mathrm{U}}(t)
    =
    |E_{\mathrm{U}}(t)|.
\end{equation}

For generating at most $K_{\max}$ candidate routes, a conservative upper bound is

\begin{equation}
    O
    \left(
    K_{\max}
    m_{\mathrm{U}}(t)
    \log n
    \right).
\end{equation}

In practice, the computational cost is lower because candidate generation is restricted by compact routing tables, anchor structures, tracked long-range links, and maximum-hop constraints.

If every ESP maintains at most

\begin{equation}
    T_{\max}
    =
    O
    \left(
    \sqrt{n}\log n
    \right)
\end{equation}

routing entries, the local search space remains significantly smaller than unrestricted global path enumeration.

\subsection{Temporal-Feasibility Filtering Complexity}

For each candidate route, the temporal-feasibility module evaluates every artificial e-link contained in the route.

Since a candidate route contains at most $H_{\max}$ artificial e-links, the availability, fidelity, coherence-budget, and resource-availability checks require

\begin{equation}
    O
    \left(
    H_{\max}
    \right)
\end{equation}

computational complexity per candidate route.

The end-to-end fidelity estimation also requires a product operation over at most $H_{\max}$ link parameters and swapping-retention factors. Therefore, its complexity is

\begin{equation}
    O
    \left(
    H_{\max}
    \right).
\end{equation}

For at most $K_{\max}$ candidate routes, the total feasibility-filtering complexity for one request is

\begin{equation}
    O
    \left(
    K_{\max}H_{\max}
    \right).
\end{equation}

For all entanglement requests in time slot $t$, the total complexity becomes

\begin{equation}
    O
    \left(
    D_tK_{\max}H_{\max}
    \right).
\end{equation}


\subsection{Time-Dependent Route-Selection Complexity}

After feasibility filtering, let

\begin{equation}
    K_{\mathrm{TF}}
    \leq
    K_{\max}
\end{equation}

denote the number of temporally feasible candidate routes.

The complete time-dependent cost evaluation of one route requires examining at most $H_{\max}$ artificial e-links. Therefore, evaluating all feasible routes requires

\begin{equation}
    O
    \left(
    K_{\mathrm{TF}}H_{\max}
    \right)
\end{equation}

computational complexity per request.

Selecting the minimum-cost route requires

\begin{equation}
    O
    \left(
    K_{\mathrm{TF}}
    \right)
\end{equation}

additional comparisons.

Because

\begin{equation}
    K_{\mathrm{TF}}
    \leq
    K_{\max},
\end{equation}

the overall route-selection complexity is bounded by

\begin{equation}
    O
    \left(
    K_{\max}H_{\max}
    \right).
\end{equation}


\subsection{Route-Repair Complexity}

When no directly feasible route is available, the algorithm evaluates the repairable candidate-route set.

For each candidate route, the controller identifies critical artificial e-links and estimates the corresponding repair delay and regeneration cost.

Since each route contains at most $H_{\max}$ artificial e-links, the repair evaluation requires

\begin{equation}
    O
    \left(
    H_{\max}
    \right)
\end{equation}

computational complexity per candidate route.

Therefore, the total repair-evaluation complexity per request is

\begin{equation}
    O
    \left(
    K_{\max}H_{\max}
    \right).
\end{equation}

This repair procedure is activated only when the directly feasible route set is empty, and therefore does not introduce additional overhead during normal routing operation.


\subsection{Regeneration-Priority Complexity}

The proactive regeneration module evaluates all currently critical artificial e-links.

Let

\begin{equation}
    m_{\mathrm{crit}}(t)
    =
    |\mathcal{E}_{\mathrm{crit}}(t)|
\end{equation}

denote the number of critical e-links at time $t$.

The calculation of regeneration-priority and utility scores requires constant computation per critical e-link. Therefore, the scoring complexity is

\begin{equation}
    O
    \left(
    m_{\mathrm{crit}}(t)
    \right).
\end{equation}

Selecting the highest-priority $R_{\max}$ regeneration candidates through complete sorting requires

\begin{equation}
    O
    \left(
    m_{\mathrm{crit}}(t)
    \log
    m_{\mathrm{crit}}(t)
    \right).
\end{equation}

Alternatively, maintaining a bounded priority queue with maximum size $R_{\max}$ reduces the selection complexity to

\begin{equation}
    O
    \left(
    m_{\mathrm{crit}}(t)
    \log
    R_{\max}
    \right).
\end{equation}

Since

\begin{equation}
    R_{\max}
    \ll
    m_{\mathrm{crit}}(t),
\end{equation}

the bounded-priority-queue implementation is preferred for large-scale quantum networks.

\subsection{Overlay-Maintenance Complexity}

The retain, regenerate, replace, and discard decisions are evaluated over the currently stored artificial e-link set.

Therefore, the basic maintenance-state evaluation requires

\begin{equation}
    O
    \left(
    m_{\mathrm{A}}(t)
    \right).
\end{equation}

If replacement candidates are selected from a bounded candidate set of size

\begin{equation}
    m_{\mathrm{new}}(t),
\end{equation}

the ranking complexity becomes

\begin{equation}
    O
    \left(
    m_{\mathrm{new}}(t)
    \log
    m_{\mathrm{new}}(t)
    \right).
\end{equation}

In the proposed implementation, replacement evaluation is restricted to a bounded candidate set instead of considering all possible ESP pairs. This limitation preserves the scalability of the compact-overlay architecture.


\subsection{Total Per-Slot Complexity}

Combining temporal-state monitoring, candidate construction, feasibility filtering, route selection, and regeneration operations, the overall computational complexity in time slot $t$ is

\begin{equation}
    O
    \left(
    m_{\mathrm{A}}(t)
    +
    D_t
    K_{\max}
    m_{\mathrm{U}}(t)
    \log n
    +
    D_t
    K_{\max}
    H_{\max}
    +
    m_{\mathrm{crit}}(t)
    \log R_{\max}
    \right).
\end{equation}

When candidate routes are generated directly from pre-established compact-routing structures rather than performing a complete graph search for every request, the route-construction component is reduced.

In this case, the dominant computational complexity becomes

\begin{equation}
    O
    \left(
    m_{\mathrm{A}}(t)
    +
    D_t
    K_{\max}
    H_{\max}
    +
    m_{\mathrm{crit}}(t)
    \log R_{\max}
    \right).
\end{equation}

Since $K_{\max}$ and $H_{\max}$ are bounded design parameters, the request-processing complexity increases approximately linearly with the number of active entanglement requests.


\subsection{Space Complexity}

The artificial-overlay graph requires

\begin{equation}
    O
    \left(
    n
    +
    m_{\mathrm{A}}(t)
    \right)
\end{equation}

storage space.

The temporal-state information maintained for each artificial e-link requires additional

\begin{equation}
    O
    \left(
    m_{\mathrm{A}}(t)
    \right)
\end{equation}

storage.

The temporary storage required for candidate routes of all active requests is bounded by

\begin{equation}
    O
    \left(
    D_t
    K_{\max}
    H_{\max}
    \right).
\end{equation}

Therefore, the overall space complexity is

\begin{equation}
    O
    \left(
    n
    +
    m_{\mathrm{A}}(t)
    +
    D_t
    K_{\max}
    H_{\max}
    \right).
\end{equation}

The compact-routing-table constraint prevents each ESP from maintaining complete global routing information. Instead of requiring $O(n)$ routing state, the proposed architecture maintains the target per-ESP routing-state complexity of

\begin{equation}
    O
    \left(
    \sqrt{n}
    \log n
    \right).
\end{equation}

Therefore, the proposed algorithm sacrifices global optimality in exchange for bounded candidate evaluation, compact routing-state requirements, and scalable temporal-state processing.

\end{enumerate}

\section{Chapter Summary}

This chapter presented the proposed temporal-state-aware compact routing and overlay-maintenance framework for quantum-native entanglement networks.

The proposed method extends compact quantum routing by incorporating the temporal evolution of stored artificial e-links into routing decisions. The main objectives were to improve temporal reliability, preserve entanglement fidelity, maintain compact routing-state requirements, control regeneration overhead, reduce request rejection, and achieve scalable decision making.

The proposed framework consists of two interacting processes: reactive routing and proactive overlay maintenance. Reactive routing handles incoming entanglement requests by constructing a bounded candidate-route set, evaluating temporal feasibility, and selecting the minimum-cost feasible route. Proactive maintenance continuously evaluates stored artificial e-links and determines whether they should be retained, regenerated, replaced, or discarded.

A temporal-state monitoring mechanism was introduced to update e-link age, effective fidelity, predicted consumption fidelity, remaining coherence budget, utilization history, and temporal classification. Based on these factors, artificial e-links are classified into safe, warning, critical, and expired states.

The candidate-route construction mechanism preserves compact routing by limiting route discovery to local e-neighbors, anchor information, and tracked long-range links. The number of candidate routes and maximum route length are explicitly bounded to avoid unrestricted path exploration.

The temporal-feasibility filtering module evaluates candidate routes using link availability, predicted fidelity, coherence budget, end-to-end fidelity, completion delay, and e-bit resource availability. This process separates directly feasible routes from repairable and infeasible routes.

A time-dependent route-selection metric was developed by combining temporal e-link quality, fidelity degradation, age, coherence risk, regeneration cost, swapping reliability, and route delay. Additional reliability-aware extensions were introduced using predicted route-success probability.

The proactive regeneration mechanism prioritizes critical artificial e-links according to temporal risk, structural importance, utilization history, predicted demand, and regeneration overhead. Regeneration decisions are constrained by quantum-memory capacity, physical-generation capability, and per-slot regeneration limits.

The overlay-maintenance strategy provides a controlled mechanism for preserving useful links, refreshing degraded links, replacing inefficient links, and removing obsolete links while maintaining compact routing constraints.

The complete routing procedure was summarized as a model-based heuristic algorithm. It does not attempt to solve the globally optimal joint routing and maintenance problem at every time slot. Instead, it achieves practical scalability through bounded candidate generation, temporal filtering, and limited regeneration operations.

Finally, the computational complexity analysis demonstrated that the proposed framework depends mainly on the number of stored artificial e-links, active requests, bounded candidate routes, maximum route length, and critical-link population. With fixed candidate and route-length limits, the request-processing complexity scales approximately linearly with the number of active requests.

The next chapter presents the simulation environment, implementation details, parameter configuration, baseline algorithms, evaluation metrics, and experimental scenarios used to validate the proposed method.